\chapter*{Preface}
\addcontentsline{toc}{chapter}{Preface}
\markboth{Preface}{Preface}

Linear algebra is the first booklet in \textit{Mathesis} that
feels like ``real'' mathematics in the sense most people mean it ---
vectors, matrices, transformations. But the point here is the structure
underneath the computation: vector spaces and bases, linear maps,
eigenvalues and diagonalization, inner product spaces, the spectral
theorem, and singular value decomposition.

It depends on Booklet~02 for the set-theoretic language of spaces and
maps, and it's a direct prerequisite for two things later in the
project: functional analysis, which is what happens when the spaces
stop being finite-dimensional, and differential geometry, which needs
tangent spaces to be vector spaces before it can put curvature on them.
